\documentclass[10pt,a4paper]{amsart}
\usepackage{amsmath,amssymb,amsthm,graphicx}
\usepackage{verbatim}
\newenvironment{reference}{\comment}{\endcomment}
%\newenvironment{reference}{}{}
\newenvironment{slogan}{\comment}{\endcomment}
\newenvironment{history}{\comment}{\endcomment}

% For commutative diagrams we use Xy-pic
\usepackage[all]{xy}

% We use 2cell for 2-commutative diagrams.
\xyoption{2cell}
% Theorem environments.
%
\theoremstyle{plain}
\newtheorem{theorem}[subsection]{Theorem}
\newtheorem{proposition}[subsection]{Proposition}
\newtheorem{lemma}[subsection]{Lemma}

\theoremstyle{definition}
\newtheorem{definition}[subsection]{Definition}
\newtheorem{example}[subsection]{Example}
\newtheorem{exercise}[subsection]{Exercise}
\newtheorem{situation}[subsection]{Situation}

\theoremstyle{remark}
\newtheorem{remark}[subsection]{Remark}
\newtheorem{remarks}[subsection]{Remarks}

\numberwithin{equation}{subsection}

% Macros
%
\def\lim{\mathop{\mathrm{lim}}\nolimits}
\def\colim{\mathop{\mathrm{colim}}\nolimits}
\def\Spec{\mathop{\mathrm{Spec}}}
\def\Hom{\mathop{\mathrm{Hom}}\nolimits}
\def\Ext{\mathop{\mathrm{Ext}}\nolimits}
\def\SheafHom{\mathop{\mathcal{H}\!\mathit{om}}\nolimits}
\def\SheafExt{\mathop{\mathcal{E}\!\mathit{xt}}\nolimits}
\def\Sch{\mathit{Sch}}
\def\Mor{\mathop{\mathrm{Mor}}\nolimits}
\def\Ob{\mathop{\mathrm{Ob}}\nolimits}
\def\Sh{\mathop{\mathit{Sh}}\nolimits}
\def\NL{\mathop{N\!L}\nolimits}
\def\CH{\mathop{\mathrm{CH}}\nolimits}
\def\proetale{{pro\text{-}\acute{e}tale}}
\def\etale{{\acute{e}tale}}
\def\QCoh{\mathit{QCoh}}
\def\Ker{\mathop{\mathrm{Ker}}}
\def\Im{\mathop{\mathrm{Im}}}
\def\Coker{\mathop{\mathrm{Coker}}}
\def\Coim{\mathop{\mathrm{Coim}}}

% Boxtimes
%
\DeclareMathSymbol{\boxtimes}{\mathbin}{AMSa}{"02}

%
% Macros for moduli stacks/spaces
%
\def\QCohstack{\mathcal{QC}\!\mathit{oh}}
\def\Cohstack{\mathcal{C}\!\mathit{oh}}
\def\Spacesstack{\mathcal{S}\!\mathit{paces}}
\def\Quotfunctor{\mathrm{Quot}}
\def\Hilbfunctor{\mathrm{Hilb}}
\def\Curvesstack{\mathcal{C}\!\mathit{urves}}
\def\Polarizedstack{\mathcal{P}\!\mathit{olarized}}
\def\Complexesstack{\mathcal{C}\!\mathit{omplexes}}
% \Pic is the operator that assigns to X its picard group, usage \Pic(X)
% \Picardstack_{X/B} denotes the Picard stack of X over B
% \Picardfunctor_{X/B} denotes the Picard functor of X over B
\def\Pic{\mathop{\mathrm{Pic}}\nolimits}
\def\Picardstack{\mathcal{P}\!\mathit{ic}}
\def\Picardfunctor{\mathrm{Pic}}
\def\Deformationcategory{\mathcal{D}\!\mathit{ef}}
\usepackage{cleveref}
\crefname{originallemma}{original lemma}{original lemmas}
\Crefname{originallemma}{Original lemma}{Original lemmas}
\makeatletter
\newlabel{coherent-lemma-acc-coherent}{{01Y8}{source}{}{sourcepointer-1}{}}
\newlabel{coherent-lemma-acc-coherent@cref}{{[originallemma][0][]01Y8}{[1][0][]source}{}{}{}}
\makeatother
\crefname{originallemma}{original lemma}{original lemmas}
\Crefname{originallemma}{Original lemma}{Original lemmas}
\makeatletter
\newlabel{divisors-lemma-universal-property-blowing-up}{{0806}{source}{}{sourcepointer-2}{}}
\newlabel{divisors-lemma-universal-property-blowing-up@cref}{{[originallemma][0][]0806}{[1][0][]source}{}{}{}}
\makeatother
\crefname{originalreference}{original reference}{original references}
\Crefname{originalreference}{Original reference}{Original references}
\makeatletter
\newlabel{lemma-regular-alteration-implies}{{0BGG}{source}{}{sourcepointer-3}{}}
\newlabel{lemma-regular-alteration-implies@cref}{{[originalreference][0][]0BGG}{[1][0][]source}{}{}{}}
\makeatother
\crefname{originallemma}{original lemma}{original lemmas}
\Crefname{originallemma}{Original lemma}{Original lemmas}
\makeatletter
\newlabel{varieties-lemma-finite-in-codim-1}{{0AB7}{source}{}{sourcepointer-4}{}}
\newlabel{varieties-lemma-finite-in-codim-1@cref}{{[originallemma][0][]0AB7}{[1][0][]source}{}{}{}}
\makeatother
\title{Characterizations of resolvability for a one dimensional integral Noetherian scheme}
\author{The Stacks Project Authors}
\date{stacks-project:resolve.tex:local-snapshot-20261012; GNU FDL 1.2 or later}
\begin{document}
\maketitle
\noindent\textit{Editorial note.} All delivered mathematical text is the author's own, copied from the pinned Stacks Project chapter file resolve.tex; no statement or proof step was rewritten, reordered or removed.\par\medskip
\noindent\textit{Editorial note.} The counted proof is the author's complete argument: the reduction of the list of implications to the construction of the sequence of blowups, the factorization of a finite resolution through blowups of closed points with the associated ascending chain of pushforwards, and the termination through the ascending chain condition together with the universal property of blowing up.\par\medskip
\noindent\textit{Editorial note.} The label lemma-resolve-curve names the counted statement, and the proof environment immediately following it in the source is delivered verbatim as its proof.\par\medskip
\noindent\textit{Editorial note.} The counted statement is carried by the author's own lemma environment in the chapter source.\par\medskip
\noindent\textit{Editorial note.} Pre-open gate: python operations/agent-authoring/tools/precard\_gate.py --source resolve.tex --label lemma-resolve-curve --start 4247 --end 4304 --local-id 28089 returned RESULT=OK (L28089 free, occupancy FREE, public interval clear). Pre-assembly tag self-check: canonical tag resolved and not present in reference/public-items.jsonl.\par\medskip
\noindent\textit{Editorial note.} The statement and proof are copied from the pinned chapter file; the cited inputs are kept as external pointers with their canonical Stacks tags.\par\medskip
\noindent\textit{Allowed standard prerequisites (editorial).}\begin{itemize}
\item Finiteness of a resolution in dimension one and of the normalization for alterations. Varieties, the lemma that a modification of a one dimensional scheme is finite, and Same chapter, the lemma deducing finiteness of the normalization from a regular alteration. Supply the finiteness of the resolution used to build the ascending chain and the implication from a regular alteration to finiteness of the normalization.
\item Ascending chain condition for coherent submodules of a coherent module. Cohomology of Schemes, the lemma that a coherent sheaf on a Noetherian scheme has the ascending chain condition on coherent submodules. Supplies the stabilization of the ascending chain of pushforwards of structure sheaves.
\item Universal property of blowing up. Divisors, the lemma giving the universal property of the blowing up of an ideal, applied to the pullback of a maximal ideal which is invertible on the resolution. Supplies the factorization of the resolution through each blowup in the constructed sequence.
\end{itemize}\medskip
\ifdefined\hypertarget\hypertarget{sourcepointer-1}{}\fi
\noindent\textit{Source reference.} Original lemma 01Y8: sources/stacks-chow/tags.txt, source lines 2547--2547. Prerequisite pointer invoked at proof line(s) 4290; it names a standard fact of a sibling chapter that the counted argument uses, and every step of the counted proof built on it is retained verbatim in the delivered body.\par\medskip
\ifdefined\hypertarget\hypertarget{sourcepointer-2}{}\fi
\noindent\textit{Source reference.} Original lemma 0806: sources/stacks-chow/tags.txt, source lines 10171--10171. Prerequisite pointer invoked at proof line(s) 4299; it names a standard fact of a sibling chapter that the counted argument uses, and every step of the counted proof built on it is retained verbatim in the delivered body.\par\medskip
\ifdefined\hypertarget\hypertarget{sourcepointer-3}{}\fi
\noindent\textit{Source reference.} Original reference 0BGG: sources/stacks-resolve/pinned-original-source.tex, source lines 3747--3747. Prerequisite pointer invoked at proof line(s) 4264; it names a standard fact of the same chapter that the counted argument uses, and every step of the counted proof built on it is retained verbatim in the delivered body.\par\medskip
\ifdefined\hypertarget\hypertarget{sourcepointer-4}{}\fi
\noindent\textit{Source reference.} Original lemma 0AB7: sources/stacks-chow/tags.txt, source lines 13010--13010. Prerequisite pointer invoked at proof line(s) 4273; it names a standard fact of a sibling chapter that the counted argument uses, and every step of the counted proof built on it is retained verbatim in the delivered body.\par\medskip
\begin{lemma}
\label{lemma-resolve-curve}
Let $Y$ be a one dimensional integral Noetherian scheme.
The following are equivalent
\begin{enumerate}
\item there exists an alteration $X \to Y$ with $X$ regular,
\item there exists a resolution of singularities of $Y$,
\item there exists a finite sequence
$Y_n \to Y_{n - 1} \to \ldots \to Y_1 \to Y$ of blowups
in closed points with $Y_n$ regular, and
\item the normalization $Y^\nu \to Y$ is finite.
\end{enumerate}
\end{lemma}

\begin{proof}
The implications (3) $\Rightarrow$ (2) $\Rightarrow$ (1) are immediate.
The implication (1) $\Rightarrow$ (4) follows from
Lemma \ref{lemma-regular-alteration-implies}.
Observe that a normal one dimensional scheme is regular hence
the implication (4) $\Rightarrow$ (2) is clear as well.
Thus it remains to show that the equivalent conditions (1), (2), and
(4) imply (3).

\medskip\noindent
Let $f : X \to Y$ be a resolution of singularities. Since the dimension
of $Y$ is one we see that $f$ is finite by
Varieties, Lemma \ref{varieties-lemma-finite-in-codim-1}.
We will construct factorizations
$$
X \to \ldots \to Y_2 \to Y_1 \to Y
$$
where $Y_i \to Y_{i - 1}$ is a blowing up of a closed point and not
an isomorphism as long as $Y_{i - 1}$ is not regular.
Each of these morphisms will be finite (by the same
reason as above) and we will get a corresponding system
$$
f_*\mathcal{O}_X \supset \ldots \supset
f_{2, *}\mathcal{O}_{Y_2} \supset
f_{1, *}\mathcal{O}_{Y_1} \supset \mathcal{O}_Y
$$
where $f_i : Y_i \to Y$ is the structure morphism.
Since $Y$ is Noetherian, this increasing sequence of coherent submodules
must stabilize
(Cohomology of Schemes, Lemma \ref{coherent-lemma-acc-coherent})
which proves that for some $n$ the scheme $Y_n$ is regular
as desired. To construct $Y_i$ given $Y_{i - 1}$ we pick a singular
closed point $y_{i - 1} \in Y_{i - 1}$ and we let $Y_i \to Y_{i - 1}$
be the corresponding blowup. Since $X$ is regular of dimension $1$
(and hence the local rings at closed points are discrete valuation
rings and in particular PIDs), the ideal sheaf
$\mathfrak m_{y_{i - 1}} \cdot \mathcal{O}_X$ is invertible.
By the universal property of blowing up (Divisors, Lemma
\ref{divisors-lemma-universal-property-blowing-up})
this gives us a factorization
$X \to Y_i$.
Finally, $Y_i \to Y_{i - 1}$ is not an isomorphism as
$\mathfrak m_{y_{i - 1}}$ is not an invertible ideal.
\end{proof}

\end{document}
